\section{Conclusion}

We built MovieMind, a reproducible movie recommendation pipeline on MovieLens 1M with explicit attention to
data sparsity, long-tail popularity, and leakage-safe evaluation.
Engineered tabular features combined with gradient boosting outperform linear models, random forests, and
our NCF implementations on test RMSE under a chronological 70/10/20 split.

\textbf{Key findings:}
\begin{itemize}
  \item More than 95\% matrix sparsity and long-tail item popularity justify hybrid features beyond pure
        collaborative filtering.
  \item Raw gradient boosting is the strongest single offline model (RMSE 0.8981, MAE 0.7062).
  \item ML and DL hyperparameter tuning improved some configurations but did not displace the raw GB
        champion in this run.
\end{itemize}

\textbf{Limitations:}
\begin{itemize}
  \item Rating prediction RMSE does not fully capture ranking quality for top-$K$ recommendation lists.
  \item Cold-start users and tail movies remain challenging despite aggregate features.
  \item DL models may benefit from longer training and richer text or genre embeddings.
\end{itemize}

\textbf{Future work:}
\begin{itemize}
  \item Two-tower retrieval and learning-to-rank objectives (Precision@K / NDCG).
  \item Session-aware or sequential recommenders using timestamps.
  \item Integration with the deployed FastAPI/Streamlit app and optional Ollama tool agent for interactive
        queries (documented separately in the repository).
\end{itemize}

The complete training and evaluation workflow is available in
\texttt{notebooks/MovieMind\_capstone.ipynb} on the project repository.
